\documentclass[11pt,a4paper]{article}

\usepackage[utf8]{inputenc}
\usepackage[T1]{fontenc}
\usepackage{geometry}
\usepackage{xcolor}
\usepackage{listings}
\usepackage{booktabs}
\usepackage{tabularx}
\usepackage{enumitem}
\usepackage{hyperref}

\geometry{margin=25mm}
\hypersetup{
  colorlinks=true,
  linkcolor=blue,
  urlcolor=blue,
  pdftitle={MolSpinPy Manual},
  pdfauthor={MolSpinPy contributors}
}

\definecolor{codebg}{RGB}{247,248,250}
\definecolor{codeframe}{RGB}{210,215,220}
\lstset{
  language=Python,
  basicstyle=\ttfamily\small,
  keywordstyle=\color{blue!65!black},
  commentstyle=\color{green!40!black},
  stringstyle=\color{red!55!black},
  backgroundcolor=\color{codebg},
  frame=single,
  rulecolor=\color{codeframe},
  breaklines=true,
  breakatwhitespace=false,
  showstringspaces=false,
  keepspaces=true,
  columns=fullflexible,
  tabsize=4
}

\title{\textbf{MolSpinPy}\\Manual for version 0.1.0}
\author{Python interface for MolSpin}
\date{Version 0.1.0}

\begin{document}
\maketitle
\tableofcontents
\newpage

\section{About this manual}

MolSpinPy provides Python objects and utilities for preparing MolSpin input,
generating parameter sweeps, running calculations, deploying work to supported
Slurm clusters, and exploring result data. This manual describes the
implementation identified as version 0.1.0 in the project package metadata.
It is based on the code currently in the repository; the release should be
treated as an early, evolving interface.

MolSpinPy prepares input and coordinates execution; it does not replace the
MolSpin simulation program. A usable MolSpin executable or a configured remote
installation is needed to perform simulations.

\section{Installation and compatibility}

The package metadata declares Python 3.8 or newer and the following runtime
dependencies: NumPy, Matplotlib, SciPy, and Paramiko. From the repository root,
install the package with:

\begin{lstlisting}
python -m pip install .
\end{lstlisting}

For development tools, install the optional development dependencies with:

\begin{lstlisting}
python -m pip install ".[dev]"
\end{lstlisting}

\subsection{Packaging notes for 0.1.0}

There are differences between the declared installation requirements and
imports in the current source:
\begin{itemize}[nosep]
  \item \texttt{MolSpinPy.processing.EigenValueSpectrumTool} imports pandas,
        but pandas is not listed in the package dependencies.
  \item The package initialization imports \texttt{typing\_extensions}, which
        is also not listed as a direct dependency.
  \item The setuptools configuration explicitly lists only the top-level
        \texttt{MolSpinPy} package, not its subpackages. Check the installed
        distribution contains \texttt{objects}, \texttt{runtime}, and the
        other subpackages; an editable install or running from the repository
        root may be needed with this packaging configuration.
  \item The source contains a Python 3.10-style union annotation in the
        rotation actions. Although the package metadata says Python 3.8+,
        Python 3.10 or newer is the practical choice for this version.
\end{itemize}
If importing the package fails because of a missing module, install the
corresponding dependency in the active environment. For example:

\begin{lstlisting}
python -m pip install pandas typing_extensions
\end{lstlisting}

\section{Core concepts}

\begin{description}[style=nextline]
  \item[MSD input] A MolSpin data file. The parser converts it into nested
        Python dictionaries that can be inspected or modified.
  \item[Project] Owns the base input, the selected MSD template, optional
        object-built spin systems, settings, tasks, an action manager, and
        local or remote execution state.
  \item[Action] Couples a target variable to an action function and optional
        timing, sequence, and parent-child settings.
  \item[ActionManager] Registers, orders, tests, and applies actions at each
        sweep step.
  \item[Template] Turns a dictionary snapshot into a MolSpin \texttt{.msd}
        input file. The default is \texttt{BasicTemplate.msd}.
  \item[Project metadata] JSON and Markdown records of the input, template,
        action setup, environment, and execution information.
\end{description}

A typical parameter-sweep workflow is:
\begin{enumerate}[nosep]
  \item Read an MSD input or assemble a project with Python objects.
  \item Create an \texttt{ActionManager} and add actions for variables to
        change.
  \item Render a new MSD input for each sweep step.
  \item Run the generated input with MolSpin locally or submit the sweep to
        a configured HPC cluster.
  \item Process result files and plot selected quantities.
\end{enumerate}

\section{Reading and rendering MSD input}

\subsection{Load an existing input file}

\texttt{ConvertMSDtoDictionary} parses MolSpin blocks into a dictionary with
top-level sections for \texttt{spinsystems}, \texttt{settings},
\texttt{actions}, \texttt{outputs}, and \texttt{tasks}. The parser accepts
line comments beginning with \texttt{//}, braces for blocks, and semicolon
terminated assignments. It tracks original case while normalizing lookup
keys for parsing.

\begin{lstlisting}
import MolSpinPy as msp

input_data = msp.msd.MsdHandling.ConvertMSDtoDictionary(
    "base_inputfile.msd"
)
if input_data is None:
    raise FileNotFoundError("Could not read the MSD input file")

project = msp.runtime.Project("example")
project.attach_dict(input_data)
\end{lstlisting}

An MSD file can also be attached directly with
\texttt{project.attach\_MSD\_file(filename)}. To add a dictionary to an
already populated project, \texttt{attach\_dict} accepts
\texttt{dictMode.CONCATINATE} from
\texttt{MolSpinPy.runtime.Project}; otherwise the default mode replaces the
base dictionary. The project validates input structure against its template
before attaching it.

\subsection{The default input shape}

The built-in template expects the following principal sections:

\begin{lstlisting}
{
    "spinsystems": {
        "system_name": {
            "spins": {},
            "interactions": {},
            "states": {},
            "transitions": {},
            "properties": {},
            "pulses": {},
            "operators": {}
        }
    },
    "settings": {},
    "actions": {},
    "outputs": {},
    "tasks": {}
}
\end{lstlisting}

The template renders spin systems, spins, interactions, states, transitions,
properties, pulses, operators, settings, outputs, and run tasks when those
sections are present. The alternative \texttt{BasicTemplate\_withActions.msd}
also renders actions and pulse sequences. MolSpin-specific block names,
property names, units, and valid task types must follow the MolSpin program's
input requirements.

\subsection{Render a single input file}

The template can be used independently of the project deployment pipeline.
This is useful for inspecting generated input or validating the rendered
configuration before a large run.

\begin{lstlisting}
from pathlib import Path

template = msp.msd.FileTemplate.get_template("BasicTemplate.msd")
rendered = template.render(input_data)
Path("preview.msd").write_text(rendered, encoding="utf-8")
\end{lstlisting}

Custom templates can be loaded with
\texttt{msp.msd.FileTemplate.get\_template(filename, folder=...)}. Template
syntax supports variable substitutions, \texttt{if}/\texttt{elif}/
\texttt{else} conditions, and loops over dictionary entries. Template
resources are packaged under \texttt{MolSpinPy/resources/templates}.

\section{Building models with Python objects}

The object API is available under \texttt{msp.objects}. A spin system can
contain electron and nuclear spins, interactions, states, transitions, and
properties. The supported interaction classes include Zeeman, zero-field
splitting, hyperfine, dipolar, and exchange interactions.

\begin{lstlisting}
from MolSpinPy.maths.vector import Vector

system = msp.objects.spin_system.SpinSystem("GS")
electron = system.add_electron(
    msp.objects.spin_object.Electron("e1")
)
nucleus = system.add_nucleus(
    msp.objects.spin_object.Nucleus("n1")
)

system.add_interaction(
    msp.objects.interactions.Zeeman_Interaction(
        "field", electron, Vector("0 0 1")
    )
)
system.add_interaction(
    msp.objects.interactions.Hyperfine_Interaction(
        "hf", electron, nucleus,
        msp.maths.tensor.Isotropic("1.0")
    )
)
system.add_state(msp.objects.State("initial", {electron: 1, nucleus: 1}))
system.set_property("initialstate", "initial")

project = msp.runtime.Project("object_model")
project.add_spin_system(system)
project.set_settings(msp.objects.Settings("general", steps=20))
\end{lstlisting}

The spin indices passed to \texttt{State} are resolved using each spin's
index mapping. A state may also be supplied as raw MolSpin text using
\texttt{State.from\_text}. Nuclei can cause state definitions to be expanded
into nuclear permutations when states are added to a spin system.

\texttt{Settings(name, steps=1, **properties)} creates a general settings
block. \texttt{Task(name, type, **properties)} creates a run task; the type
and properties must be valid for the MolSpin executable. Use
\texttt{project.add\_task(task)} to include it in the project.

\section{Actions and parameter sweeps}

\subsection{Binding an action}

An action can target a dictionary variable with a dotted key path, or a
property on a Python object. Dictionary-bound action keys are resolved
case-insensitively. The following example is based on the checked-in
nitrogen-vacancy strain-relaxation example and expects the corresponding MSD
file to contain the named interaction properties.

\begin{lstlisting}
import math
import MolSpinPy as msp

input_data = msp.msd.MsdHandling.ConvertMSDtoDictionary(
    "NV_centre_N14_strain_relaxation.msd"
)
if input_data is None:
    raise FileNotFoundError("Could not read the NV MSD input")

manager = msp.runtime.ActionManager()
project = msp.runtime.Project("NV_sweep")
project.attach_dict(input_data)
project.attach_ActionManager(manager)

shear = 1 / math.sqrt(2)
direction = msp.maths.list.List(
    f"1 1 1 {shear:.6f} {shear:.6f} {shear:.6f}"
)

function = msp.objects.Actions.LogSpaceList(
    start=1e-5,
    stop=1e-3,
    num=100,
    direction=direction,
    Normalize=False,
    Exponent=False,
    length=6,
)
action = msp.objects.Action(
    input_data,
    function,
    duration=100,
    loop=True,
    key="spinsystems.GS.interactions.strain.E",
)
manager.add_action(action)
\end{lstlisting}

The \texttt{Action} constructor accepts timing options including
\texttt{duration}, \texttt{start}, \texttt{interval}, \texttt{loop}, and
\texttt{active}; an optional \texttt{parent} creates a dependent action.
Actions are tested when first run and then applied in priority order. Action
functions update the bound value on each applicable step.

\subsection{Built-in action functions}

The built-in action package currently exports:
\begin{itemize}[nosep]
  \item \texttt{LineSpaceScalar} and \texttt{LineSpaceVector};
  \item \texttt{LogSpaceScalar}, \texttt{LogSpaceVector}, and
        \texttt{LogSpaceList};
  \item \texttt{FibonacciSphere} and \texttt{RotateVector}.
\end{itemize}

The line- and log-space functions generate numeric sequences. Vector and
list variants multiply the sequence by a direction; their normalization
options default to true where supported. The log-space functions accept
\texttt{Exponent=False} when \texttt{start} and \texttt{stop} are values
rather than base exponents. \texttt{RotateVector} rotates a three-dimensional
bound value using a rotation axis and angle.

\subsection{Custom action scripts}

Decorate a Python function with \texttt{msp.objects.ActionScript}. It receives
the current local step, bound variable, and action state when those arguments
are present in its signature. It must return a replacement value or modify
\texttt{var.value}.

\begin{lstlisting}
@msp.objects.ActionScript
def scale_value(step, var, state, factor=2.0):
    return var.value * factor

scale_action = msp.objects.Action(
    input_data,
    scale_value,
    key="spinsystems.GS.interactions.strain.E",
    factor=2.0,
)
manager.add_action(scale_action)
\end{lstlisting}

The decorated function's extra arguments become action parameters. An
\texttt{ActionManager} can also attach actions to object-bound variables,
for example \texttt{electron.spin} or a supported interaction property.
Object-bound actions are recorded only when the target object belongs to a
spin system attached to the project.

\section{Local generation and execution}

\texttt{Project.Deploy()} is the intended combined workflow for generating
per-step MSD files, writing metadata, and optionally running a local MolSpin
binary. The binary can be found through the \texttt{molspin} argument,
\texttt{MOLSPIN\_PATH} or \texttt{MOLSPIN\_HOME}, the system path, or common
checkout folders. Additional options include \texttt{steps},
\texttt{max\_processes}, and \texttt{NPZ}.

\begin{lstlisting}
project.Deploy(molspin=r"C:\path\to\molspin.exe")
\end{lstlisting}

\subsection{Output handling}

The local deployment creates \texttt{outputs/<project-name>/}. If that
project-specific directory already exists, it is removed and recreated;
directories for other project names are not removed. Preserve any previous
results for the same project name before deploying again. The
\texttt{AssertFlags} options can disable binary lookup or MolSpin execution
while retaining input generation and integrity verification. For inspecting
one rendered input independently, use \texttt{FileTemplate.render} as shown
earlier.

On a successful local run, logs are placed in
\texttt{molspin\_logs/}, task data in \texttt{molspin\_results/}, and, when
\texttt{NPZ=True}, converted archives in \texttt{molspin\_npz/}.

\section{HPC deployment}

Remote deployment uses SSH/SFTP through Paramiko and submits Slurm jobs. You
need an SSH username, host, private-key path, remote working directory, and a
MolSpin checkout or executable on the cluster. A proxy host may be supplied
when required by the cluster's SSH configuration.

\begin{lstlisting}
project.LoadHPC(
    hpc_username="user",
    hpc_host="cluster.example.org",
    hpc_key_path="~/.ssh/id_ed25519",
    remote_working_dir="/home/user/molspin-runs",
    cluster_profile="gwdg",
)

job_ids = project.SubmitToHPC(
    hpc_username="user",
    hpc_host="cluster.example.org",
    hpc_key_path="~/.ssh/id_ed25519",
    remote_binary_path="/home/user/MolSpin",
    remote_working_dir="/home/user/molspin-runs",
    partition="standard96",
    walltime="01:00:00",
    max_concurrent=20,
    max_nodes=10,
)
project.WatchHPC(job_ids)
\end{lstlisting}

\texttt{SubmitToHPC} also accepts the cluster profile and optional scheduling
settings. Built-in profiles are:
\begin{center}
\begin{tabularx}{\linewidth}{l l X}
\toprule
Profile & Default partition & Notes \\
\midrule
\texttt{gwdg} & \texttt{standard96} & Uses a manually installed Armadillo
path and a Python virtual environment. \\
\texttt{rosa} & \texttt{rosa.p} & Loads cluster modules for Armadillo and
SciPy; applies make flag overrides. \\
\bottomrule
\end{tabularx}
\end{center}

Profiles can be registered or customized with
\texttt{HpcProfile.register} and \texttt{HpcProfile.with\_overrides}.
Calling \texttt{HpcProfile.register(profile, remember=True)} also saves the
profile to \texttt{hpc\_profiles.json} inside the installed MolSpinPy package
(override the location with the \texttt{MOLSPINPY\_HPC\_PROFILES} environment
variable), so it is loaded automatically in every later script on that
machine. \texttt{HpcProfile.saved()} lists remembered profiles and
\texttt{HpcProfile.forget(name)} removes one. Built-in profiles cannot be
remembered or removed, and reinstalling or upgrading the package may wipe the
file. Use
\texttt{InstallBinary} to clone and build MolSpin at a new remote location,
or \texttt{UpdateBinary} to synchronize and rebuild an existing checkout.
With the default Git assertion enabled, updating may reset the remote checkout
to its upstream branch and remove untracked files. Review the remote checkout
before using this operation.

Remote submissions recreate the project's remote workspace directory.
Consequently, a subsequent submission for the same project name and working
directory may replace the earlier workspace. Preserve needed results before
resubmitting. Use \texttt{DownloadFromHPC} to retrieve a completed job's
results; optional email and cluster-specific module settings are also
supported by the deployment implementation. Downloading results also
removes the remote project workspace and its archive after a successful
transfer, so do not rely on that workspace as a long-term backup.

\section{Metadata and project restoration}

Metadata generation is enabled by default. Local deployment writes a JSON
record and a readable Markdown record alongside the generated files. The
record includes the base input, template source, settings, action tree,
function configuration, and execution details.

\texttt{ProjectLoader} reconstructs a project from a metadata JSON file or a
dictionary. Built-in action functions can be restored without executing
custom code. Custom \texttt{ActionScript} source is executable content and is
disabled by default:

\begin{lstlisting}
loader = msp.runtime.ProjectLoader(
    "outputs/project/project_metadata.json",
    allow_code=True,
)
restored_project = loader.build()
\end{lstlisting}

Only load metadata from a trusted source when enabling \texttt{allow\_code}.
Modules outside the built-in allow-list are refused unless explicitly
permitted with \texttt{allow\_modules}. The loader verifies recorded hashes
for custom source and dependencies; hash verification is not a substitute for
reviewing code before executing it.

\section{Processing results}

\subsection{Converting MolSpin data files}

\texttt{ProcessSingleFile(filename, output\_dir)} reads a whitespace-delimited
data file whose first row contains column names, then saves an NPZ archive
with parsed columns. This is used by local or downloaded-result workflows
when NPZ conversion is enabled.

\subsection{Time evolution plots}

\texttt{TimeEvolutionTool} accepts an NPZ file path or a dictionary of
NumPy arrays. It locates a time column and state/population arrays, then
provides line and stacked-area plots.

\begin{lstlisting}
tool = msp.processing.TimeEvolutionTool(
    "results.npz",
    time_key="time(ns)",
)
figure, axes = tool.plot_evolution()
figure.savefig("time-evolution.png", dpi=160)
\end{lstlisting}

When data are in a project workspace, a project can be passed to the tool to
resolve local or HPC result paths. The optional time key, state key(s), labels,
and cache directory can be specified when the file's column names require
explicit selection.

\subsection{Eigenvalue spectra}

\texttt{EigenValueSpectrumTool} accepts a whitespace-delimited file path or a
pandas DataFrame. It recognizes eigenvalue columns named with
\texttt{.lambdaN}, reference-state fractions named with \texttt{.refN(state)},
and pairwise mixing columns named \texttt{RefMix(a,b)} or
\texttt{EigenMix(a,b)}. Its public analysis methods include
\texttt{plot\_spectrum}, \texttt{plot\_mixing\_fractions}, and
\texttt{find\_anticrossings}. Values and thresholds use the units and scales
in the input data.

\section{Troubleshooting}

\begin{description}[style=nextline]
  \item[MSD file not found] Check the current working directory or pass an
        absolute path. The parser returns \texttt{None} for a missing file.
  \item[Invalid action key] Confirm the dot-separated path exists in the
        parsed dictionary. Use dictionary keys from the input file; key
        matching is case-insensitive.
  \item[Action type or shape error] Check the value type and dimensions of
        the action result against its target. Actions are tested by the
        manager before the first normal step.
  \item[MolSpin executable not found] Pass the executable or installation
        directory through \texttt{Deploy(molspin=...)}, set
        \texttt{MOLSPIN\_PATH}, or add MolSpin to \texttt{PATH}.
  \item[Custom action refused by the loader] Set \texttt{allow\_code=True}
        only after reviewing the metadata source, and explicitly permit any
        required external modules.
  \item[Local deployment verification fails] This is a known 0.1.0 output
        path mismatch described in the local deployment section.
\end{description}

\section{Public module map}

\begin{center}
\begin{tabularx}{\linewidth}{l X}
\toprule
Module & Main purpose \\
\midrule
\texttt{msp.msd} & MSD parsing, templates, rendering, and data conversion. \\
\texttt{msp.objects} & Spin models, actions, states, transitions, settings,
and tasks. \\
\texttt{msp.maths} & Vector, list, and tensor representations. \\
\texttt{msp.runtime} & Projects, action management, HPC profiles, and
metadata loading. \\
\texttt{msp.processing} & Time-evolution and eigenvalue-spectrum processing. \\
\texttt{msp.utility} & Runtime flags and supporting utilities. \\
\bottomrule
\end{tabularx}
\end{center}

\end{document}
